\subsection{Why: KV reuse is not stationary}
\label{sec:why}
Buffer-management theory from A0 to LRU-K assumes that each page is referenced by a stationary process.
Under that assumption two estimators need nothing but the counts: Robbins' identity predicts the mean
future count of a block seen $k$ times from the number of blocks seen $k$ and $k{+}1$ times, and
Good--Turing's missing mass predicts the share of accesses that go to blocks not yet seen. Both hold on
our traces once their requests are shuffled into random order: the ratio of actual to predicted reuse is
1.04 and 1.05 for blocks seen once, and the missing mass matches the share of accesses to new blocks within
0.002. Both fail on the traces as recorded. Blocks seen one to four times are reused
0.12--0.27 times as often as the identity predicts (0.51 and 0.65 for blocks seen five times), and new
blocks receive 82.2\% and 57.4\% of the accesses after the split where the missing mass predicts 50.1\% and
36.4\%. Three of Alibaba's four traces depart the same way (0.19--0.31 for blocks seen once); the fourth,
traceB, is nearly stationary (0.71), and it is where the estimated prior gains least. The synthetic
control departs the other way (1.95). A KV block is created with a session and stops being read when the
session ends, so its count says more about how long it has been alive than about how often it will be
read: the population is born and dies rather than being sampled.

Two consequences follow, and we tested both. If blocks die with sessions, they die quickly, and the dead
share should not depend much on how long the evaluation waits. It does not: with the split at a quarter
of the trace, cutting the window after it from about 42 minutes to 7.5 moves the dead share from 87.5\% to
90.2\% on conversation and from 90.5\% to 91.4\% on toolagent, by 1.8 and 0.2 points between 15 minutes and
the end. The window's length is nonetheless one more convention for the reference: as it lengthens, more
of the oracle's advantage falls on blocks created inside it, and the reference with dead blocks last falls
from 89.7\% to 47.5\% on conversation, while the tie rule still moves it by 38--52 points at every length.
Second, a block's age should mark it as dead better than its count. It does, but that is not enough: the
rule we registered in advance, parking blocks seen at most twice and older than the median, parks blocks
that are dead 98.0\% and 99.9\% of the time but reaches only half of the dead ones, and captures 28.0\% and
28.8\%, less than the seen-once rule. Parking the oldest three quarters, which we tried afterwards,
captures 49.0\% and 50.4\%, but on Alibaba's traces it beats the seen-once rule on two, falls below it on
the other two, and loses 19.1\% on traceB, the nearly stationary trace, where age says least; the
registered rule is below the seen-once rule on all four. Every rule we tried gains on every trace except
that one, and none beats the seen-once rule on every trace. Measuring age by the session instead of the
block, with the session taken from the prefix as the block after the first branching point of its hash
chain, changes the result by less than a point (29.8\% and 30.6\% against 30.5\% and 31.1\%), since reading
a block reads its prefix. With every rule's precision at 0.95 or more to two decimals, what decides a rule's
value is how many of the dead blocks it reaches.
